\section{Визначення параметрів зони передруйнування}
\label{sec:derivation}

\subsection{Перетворення Мелліна і функціональне рівняння}
\label{subsec:mellin}

Для побудови локального розв'язку вводимо безрозмірну радіальну
координату
\begin{equation*}
  \rho=\frac{r}{d_i}
\end{equation*}
і застосовуємо інтегральне перетворення Мелліна до крайової задачі
\eqref{eq:interface-boundary-conditions}--\eqref{eq:matched-far-field}.
Трансформанти коригувального нормального напруження поза зоною та
градієнта нормального переміщення в зоні визначаємо відповідно як
\begin{equation*}
\begin{aligned}
  \Phi_i^+(p)
    &=\int_1^\infty
      \widehat{\sigma}_i(\rho d_i)\rho^p\,\dd\rho,\\
  \Phi_i^-(p)
    &=\frac{E_i}{2(1-\nu_i^2)}
      \int_0^1
      \left.
      \frac{\partial u_\theta}{\partial r}
      \right|_{\substack{r=\rho d_i\\\theta=\beta_i}}
      \rho^p\,\dd\rho.
\end{aligned}
\end{equation*}
Оскільки $\widehat{\sigma}_i(r)=o(r^{-1})$ при $r\to\infty$,
інтеграл для $\Phi_i^+$ визначає аналітичну функцію при
$\operatorname{Re}p<0$. У вужчій області
$\operatorname{Re}p<-1-\lambda$, де збігається також трансформанта
повного напруження
\begin{equation*}
  \widetilde{\Phi}_i^+(p)
  =\int_1^\infty
    \sigma_\theta(\rho d_i,\beta_i)\rho^p\,\dd\rho,
\end{equation*}
ці дві функції пов'язані співвідношенням
\begin{equation*}
  \Phi_i^+(p)
  =\widetilde{\Phi}_i^+(p)
   +\frac{CQ_i d_i^\lambda}{p+\lambda+1},
\end{equation*}
яке явно вилучає полюс, породжений полем зрощування. Поблизу кутової точки градієнт переміщення
має порядок $r^{\lambda_i}$, де $\lambda_i$ визначено нижче, тому
$\Phi_i^-$ є аналітичною при
$\operatorname{Re}p>-1-\lambda_i$. Верхні індекси в $\Phi_i^+$ і
$\Phi_i^-$ позначають функції, віднесені відповідно до лівої та правої
півплощин, а не знак фізичної величини.

На активній ділянці $0<\rho<1$ коригувальне напруження дорівнює
$\sigma_i-CQ_i d_i^\lambda\rho^\lambda$, звідки
\begin{equation*}
  \int_0^1
  \left(\sigma_i-CQ_i d_i^\lambda\rho^\lambda\right)
  \rho^p\,\dd\rho
  =
  \frac{\sigma_i}{p+1}
  -\frac{CQ_i d_i^\lambda}{p+\lambda+1}.
\end{equation*}
Після виключення трансформант внутрішнього поля отримуємо функціональне
рівняння
\begin{equation}
\label{eq:wiener-hopf-functional}
  \Phi_i^+(p)
  +\frac{\sigma_i}{p+1}
  -\frac{CQ_i d_i^\lambda}{p+\lambda+1}
  =-\tan(\pi p)G_i(p)\Phi_i^-(p),
\end{equation}
яке спочатку справедливе у спільній смузі збіжності
\begin{equation*}
  \max\{-1-\lambda,-1-\lambda_i\}
  <\operatorname{Re}p<0.
\end{equation*}
У цій смузі обидві частини є представленнями одного аналітичного
співвідношення. Функціональне рівняння далі задає аналітичне
(за наявності явних раціональних членів --- мероморфне) продовження
необхідних комбінацій до факторизаційного контуру
$\operatorname{Re}p=0$. Сам інтеграл, що визначає $\Phi_i^+$, при цьому
не стверджується збіжним на уявній осі. Ядро рівняння має
вигляд
\begin{equation*}
  G_i(p)=\frac{D_i(p)\cos(\pi p)}{D(p)\sin(\pi p)}.
\end{equation*}
Тут $D(p)$ є характеристичним визначником задачі про кутову точку з
міжфазною зсувною тріщиною, наведеним у додатку~\ref{app:asymptotics}.
Для плоскої деформації вводимо
\begin{equation*}
  e=\frac{1+\nu_2}{1+\nu_1}\frac{E_1}{E_2},
  \qquad
  \kappa_i=3-4\nu_i.
\end{equation*}
Визначники, що відповідають зоні передруйнування в матеріалах 1 і 2,
дорівнюють
\begin{equation*}
\begin{aligned}
  D_1(p)={}&2e(1+\kappa_2)
    [\cos(2p\alpha)-\cos(2\alpha)]
    \{\sin^2[p(\pi-\alpha)]-p^2\sin^2\alpha\}\\
  &+(1+\kappa_1)
    [\sin(2p\alpha)+p\sin(2\alpha)]
    [\sin(2p(\pi-\alpha))-p\sin(2\alpha)],\\[2pt]
  D_2(p)={}&2(1+\kappa_1)
    [\cos(2p(\pi-\alpha))-\cos(2\alpha)]
    [\sin^2(p\alpha)-p^2\sin^2\alpha]\\
  &+e(1+\kappa_2)
    [\sin(2p(\pi-\alpha))-p\sin(2\alpha)]
    [\sin(2p\alpha)+p\sin(2\alpha)].
\end{aligned}
\end{equation*}
Зокрема, показник сингулярності після утворення зони в матеріалі $i$
визначається рівнянням
\begin{equation*}
  D_i(-1-\lambda_i)=0,
\end{equation*}
де $\lambda_i$ є найменшим дійсним коренем у $(-1,0)$, неперервно
пов'язаним із відповідною фізичною гілкою. Як і раніше, випадок
$\alpha=\pi/2$ розуміємо лише як вироджену границю.

\subsection{Факторизація Вінера--Гопфа}
\label{subsec:wiener-hopf}

Рівняння~\eqref{eq:wiener-hopf-functional} розв'язуємо методом
Вінера--Гопфа~\cite{noble1958wienerhopf} аналогічно до спорідненої
за постановкою задачі~\cite{KaminskyDudykPolishchuk2024}. Запишемо
факторизацію ядра у вигляді
\begin{equation*}
  \exp\left\{
    \frac{1}{2\pi\ii}
    \int_{-\ii\infty}^{\ii\infty}
      \frac{\log G_i(z)}{z-p}\,\dd z
  \right\}
  =
  \begin{cases}
    G_i^+(p), & \operatorname{Re}p<0,\\
    G_i^-(p), & \operatorname{Re}p>0,
  \end{cases}
\end{equation*}
і введемо допоміжні гамма-функціональні множники
\begin{equation*}
  K^\pm(p)=
  \frac{\Gamma(1\mp p)}{\Gamma(\tfrac12\mp p)}.
\end{equation*}
За нормування $G_i^\pm(p)\to1$ на нескінченності розв'язок
функціонального рівняння має вигляд
\begin{equation}
\label{eq:wiener-hopf-solution}
\begin{aligned}
  \Phi_i^+(p)={}&
  -\frac{pG_i^+(p)}{K^+(p)}
  \Bigg\{
    \frac{\sigma_i}{p+1}
    \left[
      \frac{K^+(p)}{pG_i^+(p)}
      +\frac{K^+(-1)}{G_i^+(-1)}
    \right]\\
  &\hspace{37mm}
    -\frac{CQ_i d_i^\lambda}{p+\lambda+1}
    \left[
      \frac{K^+(p)}{pG_i^+(p)}
      +\frac{K^+(-\lambda-1)}
      {(\lambda+1)G_i^+(-\lambda-1)}
    \right]
  \Bigg\},
  \qquad \operatorname{Re}p<0,\\[3pt]
  \Phi_i^-(p)={}&
  K^-(p)G_i^-(p)
  \Bigg\{
    \frac{\sigma_iK^+(-1)}{(p+1)G_i^+(-1)}
    -\frac{CQ_i d_i^\lambda K^+(-\lambda-1)}
    {(\lambda+1)(p+\lambda+1)G_i^+(-\lambda-1)}
  \Bigg\},
  \qquad \operatorname{Re}p>0.
\end{aligned}
\end{equation}

Для дійсного $p<0$ ядра, що використовуються далі, є додатними й
парними на уявній осі та прямують до одиниці. Тому числове обчислення
плюс-фактора можна виконувати за дійсною формулою
\begin{equation*}
  \log G_i^+(p)
  =-\frac{p}{\pi}\int_0^\infty
    \frac{\log G_i(\ii t)}{t^2+p^2}\,\dd t,
  \qquad p<0.
\end{equation*}
У числових розрахунках контролюємо додатність ядра, його дійсність та
збіжність до одиниці на кінці контуру факторизації.

\subsection{Довжина зони передруйнування}
\label{subsec:zone-length}

Із другого співвідношення~\eqref{eq:wiener-hopf-solution} при
$p\to+\infty$ випливає
\begin{equation*}
  \Phi_i^-(p)\sim\frac{1}{\sqrt{p}}
  \left[
    \frac{\sigma_iK^+(-1)}{G_i^+(-1)}
    -\frac{CQ_i d_i^\lambda K^+(-\lambda-1)}
    {(\lambda+1)G_i^+(-\lambda-1)}
  \right].
\end{equation*}
З іншого боку, застосування теореми абелевого типу
\cite{noble1958wienerhopf} до
асимптотики~\eqref{eq:zone-tip-asymptotics} дає
\begin{equation*}
  \Phi_i^-(p)\sim-\frac{k_i}{\sqrt{2pd_i}}.
\end{equation*}
Порівняння цих співвідношень дає локальний коефіцієнт інтенсивності
напружень
\begin{equation*}
  k_i=-\sqrt{2d_i}
  \left[
    \frac{\sigma_iK^+(-1)}{G_i^+(-1)}
    -\frac{CQ_i d_i^\lambda K^+(-\lambda-1)}
    {(\lambda+1)G_i^+(-\lambda-1)}
  \right].
\end{equation*}
Накладаючи умову $k_i=0$, тобто усуваючи кореневий сингулярний доданок у прийнятій асимптотиці, отримуємо рівняння для довжини зони:
\begin{equation}
\label{eq:process-zone-length}
  d_i=
  \left[
    \frac{CQ_iK^+(-\lambda-1)G_i^+(-1)}
    {\sigma_i(1+\lambda)K^+(-1)G_i^+(-\lambda-1)}
  \right]^{-1/\lambda}.
\end{equation}
Оскільки $-1<\lambda<0$, дійсність і додатність $d_i$ вимагають
$CQ_i>0$, що узгоджується з критерієм фізичної допустимості у
підрозділі~\ref{subsec:boundary-conditions}. За незмінної геометрії
формула~\eqref{eq:process-zone-length} також дає степеневий закон
$d_i\propto C^{-1/\lambda}$.

Якщо, навпаки, фіксувати геометрію, довжину міжфазної тріщини та
зовнішнє навантаження, то з тієї самої формули випливає залежність від
опору відриву матеріалу
\begin{equation*}
  d_i\propto \sigma_i^{1/\lambda}
  =\sigma_i^{-1/|\lambda|}.
\end{equation*}
Оскільки $-1<\lambda<0$, зменшення $\sigma_i$ приводить до збільшення
довжини зони передруйнування.

Використовуючи співвідношення
\begin{equation*}
  Q_i=g_1q_i l^{\lambda_0-\lambda},
\end{equation*}
де $g_1$ і $q_i$ наведено в додатку~\ref{app:asymptotics}, із
формули~\eqref{eq:process-zone-length} отримуємо
\begin{equation}
\label{eq:normalized-process-zone-length}
  \frac{d_i}{l}=
  \left[
    \frac{Cl^{\lambda_0}}{\sigma_i}
    \frac{g_1q_iK^+(-1-\lambda)G_i^+(-1)}
    {(1+\lambda)K^+(-1)G_i^+(-1-\lambda)}
  \right]^{-1/\lambda}.
\end{equation}
Це співвідношення одночасно дає змогу перевірити припущення
$d_i/l\ll1$, покладене під час формулювання локальної задачі.

\subsection{Розкриття та енергетичний параметр}
\label{subsec:zone-parameters}

Розходження берегів зони в кутовій точці визначимо як
$\delta_i=\jump{u_\theta(0,\beta_i)}$. Із означення $\Phi_i^-(p)$
випливає
\begin{equation*}
  \delta_i
  =-2
    \left.
      \left(\frac{\partial u_\theta}{\partial r}\right)^*
    \right|_{p=0,\,\theta=\beta_i}
  =-\frac{4(1-\nu_i^2)d_i}{E_i}\Phi_i^-(0).
\end{equation*}
Підставляючи розв'язок~\eqref{eq:wiener-hopf-solution} та
використовуючи~\eqref{eq:process-zone-length}, дістаємо
\begin{equation}
\label{eq:process-zone-opening}
  \delta_i=
  -\frac{4(1-\nu_i^2)\sigma_i\lambda K^+(-1)}
  {E_i\sqrt{\pi G_i(0)}(1+\lambda)G_i^+(-1)}d_i.
\end{equation}
За фіксованих властивостей матеріалу та геометрії коефіцієнт перед $d_i$
у формулі~\eqref{eq:process-zone-opening} є сталим. Тому при зміні
зовнішнього навантаження розкриття зони зростає синхронно з її довжиною,
$\delta_i\propto d_i$. Якщо ж фіксувати навантаження і змінювати
$\sigma_i$, то з урахуванням попередньої степеневої залежності
\begin{equation*}
  \delta_i\propto \sigma_i d_i
  \propto \sigma_i^{1+1/\lambda}
  =\sigma_i^{1-1/|\lambda|}.
\end{equation*}
Показник $1-1/|\lambda|$ є від'ємним, тому зменшення опору відриву
збільшує не лише довжину зони, а й її розкриття.

Роботу $W_i$ сил опору відриву на одиницю товщини тіла в межах
прийнятої моделі визначаємо рівністю~\cite{Cherepanov1986}
\begin{equation*}
  W_i=\sigma_i\int_0^{d_i}
      \jump{u_\theta(r,\beta_i)}\,\dd r
  =-\frac{4(1-\nu_i^2)\sigma_i}{E_i}d_i^2\Phi_i^-(1).
\end{equation*}
З урахуванням~\eqref{eq:wiener-hopf-solution} маємо
\begin{equation*}
  W_i=-\frac{8\lambda(1-\nu_i^2)\sigma_i^2}
  {\pi(\lambda+2)[G_i^+(-1)]^2E_i}d_i^2.
\end{equation*}
Диференціювання за довжиною зони вздовж сім'ї рівноважних станів
за фіксованих геометрії та властивостей матеріалів визначає зональний
енергетичний параметр
\begin{equation}
\label{eq:process-zone-energy-parameter}
  J_i=\frac{\dd W_i}{\dd d_i}
  =-\frac{16\lambda(1-\nu_i^2)\sigma_i^2}
  {\pi(\lambda+2)[G_i^+(-1)]^2E_i}d_i.
\end{equation}

Для числового аналізу вводимо безрозмірні параметри
\begin{equation*}
  \sigmap=\frac{Cl^{\lambda_0}}{\sigma_i},
  \qquad
  \delta_i'=\frac{E_i\delta_i}
  {4(1-\nu_i^2)\sigma_i l},
  \qquad
  J_i'=\frac{E_iJ_i}
  {4(1-\nu_i^2)\sigma_i^2l}.
\end{equation*}
Тоді формули~\eqref{eq:normalized-process-zone-length},
\eqref{eq:process-zone-opening} і
\eqref{eq:process-zone-energy-parameter} набувають форми, зручної для
числового аналізу:
\begin{equation*}
\begin{aligned}
  \frac{d_i}{l}
    &=\left[
      \sigmap\,
      \frac{g_1q_iK^+(-1-\lambda)G_i^+(-1)}
      {(1+\lambda)K^+(-1)G_i^+(-1-\lambda)}
    \right]^{-1/\lambda},\\[2pt]
  \delta_i'
    &=-\frac{\lambda K^+(-1)}
      {\sqrt{\pi G_i(0)}(1+\lambda)G_i^+(-1)}
      \frac{d_i}{l},\\[2pt]
  J_i'
    &=-\frac{4\lambda}
      {\pi(\lambda+2)[G_i^+(-1)]^2}
      \frac{d_i}{l}.
\end{aligned}
\end{equation*}
У всіх трьох співвідношеннях використовується показник $\lambda$ задачі
про кутову точку з міжфазною зсувною тріщиною, тобто корінь
$D(-1-\lambda)=0$. Показники $\lambda_1$ і $\lambda_2$, що описують
сингулярність після появи зони, не підставляються замість $\lambda$ у
формули її параметрів.

У межах прийнятої локальної моделі величину $J_i$ використовуємо як
зональний енергетичний параметр, а не як незалежно доведений глобальний
контурний інтеграл. Аналогічно до критеріїв для маломасштабних зон біля
міжфазних тріщин~\cite{KaminskyEtAl2023}, ініціювання вторинної тріщини
можна описувати енергетичною умовою
\begin{equation*}
  J_i=J_{i\mathrm{c}},
\end{equation*}
де $J_{i\mathrm{c}}$ є незалежно заданою критичною характеристикою
матеріалу в межах цього критерію. Альтернативно можна використовувати
деформаційну умову
\begin{equation*}
  \delta_i=\delta_{i\mathrm{c}},
\end{equation*}
де $\delta_{i\mathrm{c}}$ --- незалежно задане критичне розкриття.
За фіксованих геометрії та властивостей матеріалів $\delta_i$ і $J_i$
обидва пропорційні рівноважній довжині $d_i$, тому ці два критерії
описують ту саму однопараметричну сім'ю станів і є еквівалентними лише за
узгодженого вибору відповідних критичних параметрів.
